% part: intuitionistic-logic
% chapter: propositions-as-types
% section: introduction

\documentclass[../../../include/open-logic-section]{subfiles}

\begin{document}

\olfileid{int}{pty}{int}

\olsection{جان پچھان}

لامبڈا کیلکولس حسابی عمل دا اک سادہ نمونہ اے۔ ایہ متغیراں توں بنے ٹرماں اُتے کم کردا اے۔ جے $N$ تے $M$ ٹرم نیں، تاں $(NM)$ وی ٹرم اے («$N$ نوں~$M$ اُتے لایا»)۔ جے $N$ ٹرم اے، تاں $\lambd[x][N]$ ٹرم اے۔ جے $N$ وچ متغیر~$x$ اے، تاں~$x$ دیاں میل کھاندیاں موجودگیاں $\lambd[x][N]$ وچ $\lambd[x]$ عمل کار نال بندھ جاندیاں نیں تے ہُن آزاد نہیں رہندیاں۔ $(\lambd[x][N]M)$ دی شکل والے ٹرم بارے آکھدے آں کہ اوہ $\beta$-سکڑاؤ نال $\Subst{N}{M}{x}$ بن جاندا اے (یعنی $x$ دی ہر آزاد موجودگی~$N$ وچ نوں $M$ نال بدلّن دا حاصل)۔
اک ٹرم $N$، $\beta$-تبدیلی نال~$N'$ بن جاندا اے جے $N'$، $N$ وچ کسے ذیلی ٹرم دا $\beta$-سکڑاؤ کر کے ملدا اے، تے اسیں $N \redone N'$ لکھدے آں۔ $\lambd$-کیلکولس وچ اک حسابی عمل لڑی $N \redone N_1
\redone N_2 \redone \dots$ اے۔ جے ایہ لڑی کسے ٹرم $N_k$ اُتے مُکدی اے جیہدے وچ کسے ذیلی ٹرم دا $\beta$-سکڑاؤ نہیں ہو سکدا، تاں ایہنوں \emph{نارمل} یا~$N$ دا نارمل روپ آکھدے آں۔ $\lambd$-کیلکولس وچ حسابی عمل نوں ایہو جیہیاں لڑیاں وانگ سمجھو۔ حسابی عمل اُدوں رُک جاندا اے جدوں اوہدا حاصل نارمل روپ ہووے، تے نارمل روپ ای حسابی عمل دا نتیجہ اے۔ حسابی عملاں دا ایہ سادہ نمونہ ٹیورنگ-کامل نکلدا اے۔ ہر جزوی بازگشتی فنکشن $\lambd$-کیلکولس دے کسے ٹرم نال حسابیا جا سکدا اے۔

ہیسکل کری تے ولیم ہاورڈ نے قدرتی اخذ تے $\lambd$-کیلکولس وچ نیڑی سانجھ آپو اپنے طور اُتے لبھی۔
سہج سمجھ نال، ٹرم $\lambd[x][N]$ اک فنکشن اے۔ $x$ اوہدا آرگومنٹ اے، تے $N$ دَسّدا اے کہ~$x$ دِتّا ہووے تاں قدر کیویں حسابنی اے۔ ٹرم $(\lambd[x][N]M)$ دا مطلب فیر فنکشن $\lambd[x][N]$ نوں آرگومنٹ~$M$ اُتے لاؤنا اے، تے $\beta$-سکڑاؤ دَسّدا اے کہ اوہدی قدر کیویں کڈھنی اے۔ $x$ نوں $N$ وچ~$M$ نال بدلو (تے جیہڑا ٹرم ملے، اوہدی قدر کڈھدے رہو)۔
ہُن ایہناں نوں مقرر دائرے تے حاصل دی حد والے فنکشن سمجھو۔ جے $\lambd[x][N]$ دا دائرہ $A$ اے تے اوہدی حاصل دی حد~$B$ اے، تاں متغیر~$x$ نوں ایویں سمجھو کہ اوہ صرف~$A$ توں آرگومنٹ لیندا اے تے $N$،~$B$ وچوں قدراں دیندا اے۔ سو $(\lambd[x][N]M)$ دا مطلب صرف اوہدوں ہونا چاہیدا اے جے $\lambd[x][N]$، $A \to B$ والا فنکشن اے تے $M \in A$۔
جے اسیں ایہ جانکاری $\lambd$-کیلکولس وچ شامل کرنی چاہیئے، تاں سانوں \emph{نوع بندی والا} $\lambd$-کیلکولس ملدا اے۔ نوع بندی والے $\lambd$-کیلکولس وچ ہر عبارت $(\lambd[x][N]M)$ جائز نہیں، صرف اوہ جائز نیں جیہناں وچ $\lambd[x][N]$ تے $M$ دیاں «نوعاں» آپس وچ میل کھاندیاں نیں، یعنی $\lambd[x][N]$ دی نوع $A \lif B$ اے تے $M$ دی نوع~$A$ اے۔ $\lambd[x][N]$ دی نوع $A \to B$ صرف اوہدوں اے جے $x$ دی نوع $A$ تے $N$ دی نوع~$B$ اے۔ عام طور اُتے ہر عبارت $(M_1M_2)$ جائز نہیں۔ صرف اوہ ٹرم $(M_1M_2)$ جائز نیں جیہناں وچ $M_1$ دی نوع $A \to B$ تے $M_2$ دی نوع~$A$ اے۔ جے ایویں اے، تاں $(M_1M_2)$ دی نوع~$B$ اے۔

ہُن صاف اے کہ ایہ نوع بندی دے قاعدے قدرتی اخذ دے !!{derivation}s نال میل کھاندے نیں۔ اُتھے نوعاں نوں !!{formula}s ظاہر کردے نیں تے !!{derivation}s آپ $\lambd$-ٹرماں نال میل کھاندے نیں۔ مثال لئی !!a{derivation}، $!A \lif !B$ دا، ٹرم $\lambd[\typeof{x}{!A}][N]$ نال میل کھاندا اے جیہدی فنکشن نوع $!A \lif !B$ اے۔ قدرتی اخذ دے استنباطی قاعدے نوع بندی والے لامبڈا کیلکولس دے نوع بندی قاعدیاں نال میل کھاندے نیں۔ مثال لئی:
\begin{prooftree}
  \AxiomC{$\Discharge{!A}{x}$}
  \DeduceC{$!B$}
  \DischargeRule{\Intro\lif}{x}
  \UnaryInfC{$!A \lif !B$}
  \DisplayProof\bottomAlignProof
  \quad\text{نال میل کھاندا اے}\quad
  \Axiom$x: !A \fCenter N: !B$
  \UnaryInf$ \fCenter \lambd[\typeof{x}{!A}][N] : !A \lif !B$
\end{prooftree}
سجّے پاسے دے قاعدے دا مطلب اے کہ جے $x$ دی نوع~$!A$ اے تے $N$ دی نوع~$!B$ اے، تاں $\lambd[\typeof{x}{!A}][N]$ دی نوع~$!A \lif
!B$ اے۔

$\Elim\lif$ قاعدہ جوڑ والے ٹرماں دے نوع بندی قاعدے نال میل کھاندا اے، یعنی:
\[
  \AxiomC{$!A \lif !B$}
  \AxiomC{$!A$}
  \RightLabel{\Elim\lif}
  \BinaryInfC{$!B$}
  \DisplayProof
  \quad\text{نال میل کھاندا اے}\quad
  \Axiom$\fCenter M_1: !A \lif !B$
  \Axiom$\fCenter M_2: !A$
  \BinaryInf$ \fCenter (M_1M_2) : !B$
  \DisplayProof
\]
جے \Intro\lif{} قاعدے دے جھٹ پچھوں \Elim\lif{} قاعدہ آوے، تاں !!{derivation} نوں سوکھا کیتا جا سکدا اے:
\begin{gather*}
  \AxiomC{$\Discharge{!A}{x}$}
  \DeduceC{$!B$}
  \DischargeRule{\Intro\lif}{x}
  \UnaryInfC{$!A \lif !B$}
  \AxiomC{}
  \DeduceC{$!A$}
  \RightLabel{\Elim\lif}
  \BinaryInfC{$!B$}
  \DisplayProof
  \quad\redone\quad
  \AxiomC{}
  \DeduceC{$!A$}
  \DeduceC{$!B$}
  \DisplayProof
\end{gather*}
ایہ لامبڈا ٹرماں دے $\beta$-سکڑاؤ نال میل کھاندا اے:
\[
(\lambd[\typeof{x}{!A}][N]M) \quad\redone\quad
\Subst{N}{M}{x}.
\]

قاعدیاں تے نوعاں دا ایہو جیہا میل $ \land$ دے قاعدیاں تے «حاصلِ ضرب» نوعاں وچ، تے $ \lor$ دے قاعدیاں تے «جوڑ» نوعاں وچ وی اے۔

سادہ نوع بندی والے لامبڈا کیلکولس دے ٹرماں تے قدرتی اخذ دے !!{derivation}s دے ایس میل نوں «کری–ہاورڈ»، «ثبوت بطور پروگرام»، یا «قضیے بطور نوعاں» دا میل آکھدے آں۔ !!{formula}s (قضیاں) دے نوعاں نال تے ثبوتاں دے ٹرماں نال میل توں علاوہ، اسیں ایہ میل ایویں اکٹھے وکھا سکدے آں:
\begin{center}
  \begin{tabular}{c c c}
    منطق & پروگرام \\
    \hline
    قضیہ/!!{formula} & نوع \\
    ثبوت & ٹرم \\
    مفروضہ & متغیر \\
    خارج کیتا مفروضہ & بندھیا متغیر \\
    کھُلّا مفروضہ & آزاد متغیر \\
    $!A \lif !B$ & فنکشن نوع \\
    $!A \land !B$ & حاصلِ ضرب دی نوع \\
    $!A \lor !B$ & جوڑ دی نوع \\
    $\lfalse$ & خالی نوع \\
    \hline
  \end{tabular}
\end{center}

کری–ہاورڈ میل خودکار ثبوت معاوناں تے پروگراماں دیاں نوعاں جانچن والے نظاماں دی اک بنیادی گل اے۔ وجہ ایہ اے کہ کسے قضیے دے گواہ ثبوت دی جانچ (جیویں اسیں اُتّے کیتی)، اوسے گل دی جانچ اے کہ کسے پروگرام (ٹرم) دی نوع اوہی اے جیہڑی اوہدے لئی دَسّی گئی اے۔
\end{document}
